% Abstract (part of the 0.35-page front matter)
Large language models (LLMs) cannot read execution traces, which are far larger than their context windows and describe executions absent from their training data, and when they answer anyway they can report problems that are not there. We present \sysname, an assistant developed for Ciena that answers natural-language questions about Linux kernel traces. Analyzers built on Trace Compass compute measured findings from the whole trace; small LLM decisions, checked by code, map a question to the threads it concerns; and an LLM writes the answer from the findings alone. \sysname is part of Trace Compass: each answer names, and opens, the Trace Compass views that show its evidence. To measure false alarms as well as answers, we built a symptom bench: nine C programs with 44 injectable faults covering 48 performance symptoms, 56 LTTng traces with program-written answer keys, and 516 yes/no answers per system to questions asked of traces with and without each symptom. With the same answering model, \sysname reaches a balanced accuracy of 0.93 (precision 0.85, recall 0.93), stable across three runs, against 0.58 for baselines that read trace text with or without retrieved documentation, 0.69 for an LLM given TMLL's analyses, and 0.79 for a coding agent with a shell over the whole trace, which raised nearly five times as many false alarms and took 13 times as long per answer (median). We report the lessons of developing against such a bench and the problems that remain open.
